% frontmatter/abstract.tex
\begin{center}
\pagestyle{empty}
\Large \textbf{TOM TAT}
\end{center}
\addcontentsline{toc}{section}{TOM TAT}

\vspace{0.5cm}
Bao cao nay trinh bay qua trinh quy hoach, thiet ke va trien khai mo phong he thong mang Campus
ba tang (Core -- Distribution -- Access) ket hop vung phi quan su (DMZ) su dung cong cu gia lap mang
\textbf{Mininet} tren nen tang \textbf{Ubuntu Linux}. He thong duoc xay dung theo mo hinh thiet ke phan cap ba lop
(\textit{Three-Layer Hierarchical Model}) nham dam bao tinh on dinh, kha nang mo rong va quan ly tap trung.

Cac thanh phan chinh cua he thong bao gom: 
(1)~\textbf{Topology ba lop} gom Core Router, hai Distribution Router (dist1, dist2), hai Access Switch (acc1, acc2) cung cac thiet bi dau cuoi tren VLAN\,10 va VLAN\,20; 
(2)~\textbf{Vung DMZ} chua hai may chu Web (web1: 10.10.10.11, web2: 10.10.10.12) duoc bao ve boi dmz\_router va sw\_dmz; 
(3)~\textbf{NAT/PAT} cho phep cac may chu noi bo truy cap Internet qua dia chi ngoai 203.0.113.1, dong thoi cau hinh \textit{Static NAT} anh xa web1 va web2 ra dia chi cong khai; 
(4)~\textbf{ACL va Firewall} ngan chan truy cap trai phep vao cac cong nhay cam (22, 3306) tu mang ngoai;
(5)~\textbf{Can bang tai theo nguong} voi Bo can bang tai Python tu dong chuyen may chu hoat dong giua web1 va web2 khi tai vuot nguong 80\% hoac giam xuong duoi 20\%;
(6)~\textbf{Dinh tuyen dong OSPF} thong qua FRRouting (Zebra + OSPFd) dam bao hoi tu tu dong va thu hep thoi gian khi xay ra su co.

Ket qua thuc nghiem cho thay he thong van hanh on dinh: thong luong TCP dat tren 88~Mbps sau NAT+ACL
(chi giam khoang 6.1\% so voi khong bao mat), do tre RTT duoi 4~ms, co che chuyen may chu cua Load Balancer
hoat dong dung theo nguong da thiet lap, va OSPF hoi tu thanh cong giua cac node.
Day la mo hinh mang campus tieu bieu, dap ung day du yeu cau ve bao mat, hieu nang va kha nang mo rong
cho cac doanh nghiep vua va nho.